\section{Weights, primitives, and homology suspension}
\label{sec:weights}

Homology suspension supplies two inputs to the last nonzero-suspension
argument. Its kernel identifies a terminal spherical class as a
Pontryagin square. On primitive weight components, its dual identifies
a detecting functional with the restriction of a functional at the
preceding level. We establish these statements together with finite
weight projections at level zero, where group completion requires a
filtration in place of the positive-level splitting.

For a spectrum $T$, put
\[
 D_rT=(T^{\wedge r})_{h\Sigma_r},\qquad
 \mathbb P(T)=\bigvee_{r\geq0}D_rT,\qquad D_0T=S.
\]
Thus $\mathbb P(T)$ is the free commutative ring spectrum on $T$, and its
$r$th summand has \emph{weight} $r$. Throughout, extended powers, free
algebras, pushouts, and localizations are derived constructions. We work
with their natural maps and families in the relevant homotopy coherent
categories. In particular, naturality of a construction includes its
action on mapping spaces, not just on homotopy classes of individual maps.

\subsection{A structured free-algebra model}

Let $C(X)$ be the reduced free $E_\infty$ monoid on a based space $X$:
the basepoint of $X$ is identified with the monoid unit. The usual
approximation and homology theorems identify its group completion with
$QX$ and calculate its homology by Dyer--Lashof operations
\cite[Part~I, Theorems~4.1--4.2 and Corollary~4.5]{CLM1976}.
We first record the structured splitting that will specify
our weights, including their behavior in families.

\begin{proposition}\label{prop:free-ring}
There is a natural equivalence of augmented commutative ring spectra
\[
 \Sigma^\infty_+ C(X)\simeq\mathbb P(\Sigma^\infty X).
\]
The reduced suspension spectrum of $C(X)$, based at its unit, corresponds
to the sum of the positive weights. The adjoint of projection to weight
one is the canonical map
\[
 C(X)\longrightarrow QX.
\]
This map is group completion and is an equivalence when $X$ is connected.
In the connected case, the weight-one projection is consequently the
stable counit
$\Sigma^\infty QX\to\Sigma^\infty X$.
\end{proposition}

\begin{proof}
Let $F$ denote the free unbased commutative monoid in spaces. The reduced
free monoid is the homotopy pushout
\[
 C(X)\simeq F(X)\amalg_{F(*)}*.
\]
Indeed, mapping this pushout into a commutative monoid $Y$ gives the
homotopy fibre of
$\operatorname{Map}(X,Y)\to\operatorname{Map}(*,Y)=Y$
over the unit. This is exactly the mapping space of based maps from $X$
to $Y$, based at its unit. The pushout therefore has the required reduced
free universal property, including the basepoint homotopy.

The functor $\Sigma^\infty_+$ is a strong symmetric monoidal left adjoint
from cartesian spaces to spectra. On commutative algebra objects it is
again a left adjoint, and hence preserves this pushout. Write
$T=\Sigma^\infty X$. The basepoint section and augmentation give a natural
stable splitting
\[
 \Sigma^\infty_+X\simeq S\vee T.
\]
Under it, the spectrum map induced by the basepoint is the inclusion of
the first summand. Since the free commutative algebra functor sends a
direct sum to a coproduct of algebras, the pushout after suspension is
\[
 \begin{aligned}
 \Sigma^\infty_+C(X)
 &\simeq\mathbb P(S\vee T)\amalg_{\mathbb P(S)}S\\
 &\simeq
 \bigl(\mathbb P(S)\amalg\mathbb P(T)\bigr)
       \amalg_{\mathbb P(S)}S
 \simeq\mathbb P(T).
 \end{aligned}
\]
Here $\mathbb P(S)\to S$ sends its generator to the ring unit. The
augmentation of the resulting algebra sends $T$ to zero. Thus the
positive-weight summands are precisely the augmentation fibre, or
equivalently the reduced suspension spectrum at the unit. In this
description the generating spectrum map is the \emph{difference} between
the inclusion of $X$ and its basepoint section.

We identify the weight-one adjoint using the algebra projection
\[
 \mathbb P(T)\longrightarrow S\oplus T,
\]
where the target is the split square-zero algebra. For any spectrum $M$,
the fibre over the unit point of
\[
 \mathrm{GL}_1(S\oplus M)\longrightarrow\mathrm{GL}_1S
\]
identifies with $\Omega^\infty M$, via $y\mapsto1+y$, as a grouplike
commutative monoid. To see the structure as well as the underlying space,
expand a multiproduct using distributivity in
$(S\oplus M)^{\wedge t}$. The term with no $M$ factors gives the base
unit. The terms with one $M$ factor give the sum of the inputs. Every term
with two or more $M$ factors vanishes by the square-zero structure. These
formulas respect permutations and iterated products: the structure maps
on the fibre are the additive sum-over-fibres maps. They therefore give
the stated identification with the additive infinite loop space of $M$.

The adjunction from the suspension ring now gives an $E_\infty$ map
$C(X)\to\Omega^\infty T$. On the generating based space it is the usual
map $X\to\Omega^\infty\Sigma^\infty X$. Reduced freeness identifies it
with the canonical stabilization map. Equivalently, maps from its group
completion into a grouplike commutative monoid $\Omega^\infty E$ are
based maps $X\to\Omega^\infty E$, or spectrum maps
$\Sigma^\infty X\to E$. This identifies the group completion with $QX$.
For connected $X$, the reduced free monoid has trivial $\pi_0$, so it is
already grouplike and no group completion is needed.
Under $C(X)\simeq QX$, the adjoint of the weight-one projection is then
the identity of $QX$. By the suspension-spectrum adjunction this says
exactly that the projection is the stable counit.
\end{proof}

\begin{remark}
For connected $X$, the underlying stable decomposition is the Snaith
splitting \cite[Theorem~1.1]{Snaith1974}. The argument above also
establishes the augmented commutative ring-spectrum structure and based
naturality needed for multiplication and suspension below. For a
free-algebra formulation of this structured splitting, see
\cite[Theorem~2.2 and Appendix~A]{Kuhn2006}.
\end{remark}

\subsection{Suspension and the linear coproduct}

\begin{lemma}\label{lem:weight-suspension}
Under the splitting of Proposition~\ref{prop:free-ring}, the stable
adjoint suspension on reduced pieces is weight preserving. At weight
$r$ it is the natural map
\begin{equation}\label{eq:weight-suspension}
 \Sigma D_rT\longrightarrow D_r(\Sigma T),
 \qquad T=\Sigma^\infty X,
\end{equation}
formed by using the same suspension coordinate in all $r$ factors. For
$T=\Sigma^nS$, this is the Thom spectrum map induced by the inclusion
of the diagonal trivial line into the permutation representation.

Under group completion the map is the ordinary homology suspension
$\sigma_*:H_*QX\to H_{*+1}Q\Sigma X$. It commutes with the upper-index
Dyer--Lashof operations on reduced classes and kills products of two
positive-dimensional classes. If $c$ is the class of a component point
and $z$ has positive degree, then
\[
 \sigma_*(c*z)=\sigma_*z.
\]
\end{lemma}

\begin{proof}
For $v\in S^1$, the based map $x\mapsto v\wedge x$ induces
$C(X)\to C(\Sigma X)$. These maps form a pointed family, since the
basepoint of $S^1$ induces the constant map to the monoid unit. Its
adjoint is
\[
 S^1\wedge C(X)\longrightarrow C(\Sigma X).
\]
Since $\Sigma X$ is connected, the target is $Q\Sigma X$. On the
generating space $X$, adjunction identifies this family with the usual
stabilization into $\Omega Q\Sigma X\simeq QX$. Reduced freeness and
group completion therefore identify its adjoint with the usual structure
map of the infinite loop spaces.

The positive-weight equivalence in Proposition~\ref{prop:free-ring} is
a natural equivalence between the reduced functors
\[
 X\longmapsto\operatorname{fib}(\Sigma^\infty_+C(X)\to S),
 \qquad
 X\longmapsto\bigvee_{r>0}D_r(\Sigma^\infty X).
\]
Both send the one-point based space to zero. Their action on mapping
spaces therefore respects the basepoint given by factorization through
that space, including its nullhomotopy. Apply this naturality to the
pointed mapping-space family
$v\mapsto(v\wedge-)$, and then take the stable adjunction. On the
$r$th extended power the same $v$ is used in every tensor factor. This
gives \eqref{eq:weight-suspension} with no mixing of weights. For a
sphere, $D_r(\Sigma^nS)$ is the Thom spectrum of $n\rho_r$ over
$B\Sigma_r$, where $\rho_r$ is the permutation representation. The
map in question is the Thom map associated to
\[
 n\rho_r\oplus\mathbb R\longrightarrow(n+1)\rho_r,
\]
using the diagonal line $\mathbb R\subset\rho_r$.

The commutation rule with upper-index Dyer--Lashof operations is the
usual suspension rule for free infinite loop space homology
\cite[\S20, p.~481]{May1977}. The product rules can also be read directly from the
adjunction of loop multiplication: the suspension of loop concatenation
factors through the pinch map and gives the sum of the two suspensions.
Thus a cross product of two positive-dimensional classes has zero image.
For a product with a component point, the surviving term is the
suspension of the positive-dimensional factor, multiplied by the
augmentation of the point class, which is one.
\end{proof}

For every $n\geq0$ we may equip $\mathbb P(\Sigma^nS)$ with the
\emph{linear coproduct}: it is the commutative ring map whose restriction
to the generating spectrum is the sum of its inclusions in the two
factors. At $n=0$ this is auxiliary structure and is not asserted to be
the space diagonal. At positive levels it is the space diagonal, as the
following observation shows.

\begin{lemma}\label{lem:linear-diagonal}
For a suspension $X$, the space diagonal on $C(X)$ corresponds under
Proposition~\ref{prop:free-ring} to the linear coproduct on
$\mathbb P(\Sigma^\infty X)$. Its component from weight $r$ to weights
$(i,r-i)$ is the stable transfer for
\[
 \Sigma_i\times\Sigma_{r-i}\subset\Sigma_r.
\]
\end{lemma}

\begin{proof}
In terms of the reduced generator, the diagonal is the sum of its two
linear inclusions and the term induced by the reduced diagonal
$X\to X\wedge X$. The latter is null for a suspension. Because the
source ring is free, a homotopy on the generating spectrum gives the
required homotopy of commutative ring maps.

Apply $D_r$ to the linear map $T\to T\vee T$. The terms choosing the
first summand $i$ times are indexed by the transitive $\Sigma_r$-set
$\Sigma_r/(\Sigma_i\times\Sigma_{r-i})$. Equivariantly, the map into
these terms is the diagonal into a finite sum indexed by that set.
On homotopy orbits it is the stable transfer to the indicated subgroup.
This describes every component, including the two end components.
\end{proof}

\subsection{The Steenrod modules in a primitive weight}

At weight $r>0$, let $P_r(n)$ denote the subspace of
$H_*D_r(\Sigma^nS)$ primitive for the linear coproduct: all components
with $0<i<r$ vanish. Its degreewise dual is the quotient of
$H^*D_r(\Sigma^nS)$ by the sum of the images of the proper two-block
cohomology transfers. At each fixed weight these groups are of finite
type, so this duality is exact degree by degree. For $n>0$, put
\[
 B_n=\widetilde H^*(QS^n),\qquad QB_n=B_n/B_n^2.
\]
Lemma~\ref{lem:linear-diagonal} identifies the transfer quotients with
the corresponding weight parts of $QB_n$.

At weight $2^k$, restriction to the regular elementary abelian subgroup
of $\Sigma_{2^k}$ will identify this quotient with an ideal in a
polynomial invariant ring. We first describe that ring.

For $k\geq1$, write
\[
 P_k=\F_2[x_1,\ldots,x_k],\qquad |x_i|=1,
 \qquad D_k=P_k^{\GL_k(\F_2)}
       =\F_2[d_0,\ldots,d_{k-1}],
\]
where $|d_j|=2^k-2^j$. Here $D_k$ is a Dickson algebra, in contrast to
the extended-power spectrum $D_rT$. We suppress the rank $k$ on the
generators $d_j$ and set $d_k=1$. The Dickson generators are the
coefficients of the additive polynomial
\begin{equation}\label{eq:dickson-polynomial}
 \prod_{a\in\operatorname{span}\{x_1,\ldots,x_k\}}(U+a)
   =\sum_{j=0}^k d_jU^{2^j}.
\end{equation}
The invariant-ring description is Dickson's theorem
\cite{Dickson1911}. We recall the root-polynomial proof of the
binary case, also developed in Wilkerson's account
\cite[\S I]{Wilkerson1983}.

\begin{lemma}\label{lem:dickson-ring}
The coefficients \(d_0,\ldots,d_{k-1}\) in
\eqref{eq:dickson-polynomial} are algebraically independent, have
degrees \(2^k-2^j\), and generate
\(P_k^{\GL_k(\F_2)}\).
\end{lemma}

\begin{proof}
The product over a finite binary vector space is additive in \(U\).
Indeed, if \(W'=W\oplus\langle a\rangle\) and its root polynomial
for \(W\) is \(P_W\), then
\[
 P_{W'}(U)=P_W(U)P_W(U+a)
          =P_W(U)^2+P_W(a)P_W(U).
\]
Induction from \(P_{\{0\}}(U)=U\) proves additivity and the displayed
powers of \(U\). Homogeneity gives the degrees of the coefficients.
They are invariant because an invertible binary linear change of
variables permutes the roots.

Put \(B=\F_2[d_0,\ldots,d_{k-1}]\subset P_k\) and
\(L=\F_2(x_1,\ldots,x_k)\). Each \(x_i\) satisfies the monic
root polynomial, so \(P_k\) is integral over \(B\). Therefore
\(B\) has transcendence degree \(k\); its \(k\) indicated
generators are algebraically independent and \(B\) is integrally
closed. The derivative of the root polynomial is the nonzero product
\(d_0\) of all nonzero binary linear forms. Its splitting field
over \(\operatorname{Frac}B\) is \(L\). Every automorphism of
this splitting field acts invertibly and linearly on the root space;
conversely every such linear transformation fixes the coefficients.
Consequently \(L^{\GL_k(\F_2)}=\operatorname{Frac}B\).
An invariant polynomial lies in this fraction field and is integral
over \(B\), hence belongs to \(B\). This proves the assertion.
\end{proof}

\begin{lemma}\label{lem:weight-modules}
The only nonzero positive-weight transfer quotients are at weights
$r=2^k$, $k\geq0$. As Steenrod modules they are
\begin{equation}\label{eq:weight-modules}
 M_0(n)=\Sigma^n\F_2,\qquad
 M_k(n)=\Sigma^n d_0^nD_k\quad(k\geq1),\qquad n\geq0.
\end{equation}
For $n>0$ these give the weight decomposition of $QB_n$. The action on
$d_0^nD_k$ is its ordinary action as an ideal in $P_k$.

The dual of \eqref{eq:weight-suspension} on these quotients is, after
the common grading suspension, the inclusion
\begin{equation}\label{eq:module-suspension}
 d_0^{n+1}D_k\hookrightarrow d_0^nD_k
 \qquad(k\geq1).
\end{equation}
In weight one it is the identity with the corresponding shift.
\end{lemma}

\begin{proof}
We first determine dimensions and the weights in which primitives can
occur. The free infinite loop space homology theorem gives, for $n>0$,
\begin{equation}\label{eq:free-homology-basis}
 H_*QS^n=
 \F_2\bigl[Q_{j_1}\cdots Q_{j_k}g_n:
          0<j_1\leq\cdots\leq j_k\bigr],
\end{equation}
including the empty word $g_n$, where
$Q_jz=Q^{j+|z|}z$ and $|g_n|=n$ \cite{CLM1976}.
The generator $g_n$ is primitive. Naturality of the operations for the
coproduct, together with Cartan and the vanishing of positive
operations on the unit, shows that every displayed polynomial
generator is primitive.

For $n=0$, the same polynomial description holds for
$H_*\mathbb P(S)$ on the free degree-zero generator, which we denote
here by $u$. One may see this from the free unbased monoid
$\coprod_rB\Sigma_r$, whose suspension ring is free on a degree-zero
spectrum generator. The Dyer--Lashof operations for commutative ring
spectra agree with the space operations on such suspension rings
\cite{BMMS1986}. Naturality for the \emph{linear} ring coproduct again
makes the polynomial generators primitive. This argument does not
identify that coproduct with the actual space diagonal at level zero.

In both calculations, the bottom generator has weight one, a
Dyer--Lashof operation doubles weight, and multiplication adds weights.
In a polynomial Hopf algebra on primitive generators, the primitives
have basis the iterated $2$-powers of those generators. For example,
this follows by writing the primitive identity as
$f(\mathbf x+\mathbf y)=f(\mathbf x)+f(\mathbf y)$: a monomial involving
distinct variables produces a nonzero cross term, and for a single
variable all interior binomial coefficients vanish exactly when the
exponent is a power of two. Every polynomial involves only finitely
many variables, so the argument also applies to the infinite
polynomial algebra.

It follows that a primitive basis is given by the words
\[
 Q_{j_1}\cdots Q_{j_k}g_n,
 \qquad 0\leq j_1\leq\cdots\leq j_k,
\]
with $u$ in place of $g_n$ at $n=0$. Initial zero indices record
iterated squares. A word of length $k$ has weight $2^k$ and degree
\[
 2^kn+j_1+2j_2+\cdots+2^{k-1}j_k.
\]
Thus primitives vanish at the other positive weights. Setting
$a_1=j_1$ and $a_i=j_i-j_{i-1}$ for $i>1$ rewrites this degree as
\begin{equation}\label{eq:primitive-degree-count}
 2^kn+\sum_{i=1}^k(2^k-2^{i-1})a_i,
 \qquad a_i\geq0.
\end{equation}
These are exactly the monomial degrees in
$\Sigma^nd_0^nD_k$. At $n=0$ we perform this count at fixed weight,
allowing degree zero; it remains finite in each degree at that weight.

We next identify the Steenrod module, rather than just its graded
dimension. Fix $r=2^k$ with $k\geq1$ and let
$V=(\Z/2)^k\subset\Sigma_r$ be the regular elementary abelian
subgroup. The extended power $D_r(\Sigma^nS)$ is the Thom spectrum
of $n\rho_r$ over $B\Sigma_r$. Restrict its cohomology to $BV$.
Every proper two-block transfer restricts to zero. Indeed, the double
coset formula expresses its restriction as transfers from subgroups
\[
 V\cap g(\Sigma_i\times\Sigma_{r-i})g^{-1},\qquad 0<i<r.
\]
Each is proper, because $V$ acts transitively on the particles whereas
a two-block subgroup preserves a nonempty proper subset. If $W<V$,
restriction $H^*BV\to H^*BW$ is surjective. The projection formula
then expresses transfer of any class as a class from $H^*BV$ times
the transfer of $1$, which is the even index $[V:W]$, hence zero.
Naturality of the mod-$2$ Thom class gives the same assertion with
the representation twist.

The normalizer of $V$ acts through $\GL_k(\F_2)$, so the restriction
image is contained in $D_k$ times the Thom class. It contains all
of that module: the Stiefel--Whitney classes of $\rho_r$ restrict to
the Dickson generators in \eqref{eq:dickson-polynomial}. More
explicitly, $\rho_r|_V$ is the sum of all the real one-dimensional
characters, and its total Stiefel--Whitney class is the product of
$1+a$ over all $a\in V^*$. Thus
$w_{2^k-2^j}(\rho_r)|_{BV}=d_j$, and polynomials in these classes,
multiplied by the Thom class, give the claimed image. Restriction
therefore gives a surjection from the transfer quotient onto $D_k$
times the Thom class. The dimension computation
\eqref{eq:primitive-degree-count} makes this surjection an
isomorphism in every degree.

Finally, split the regular representation over $BV$ as
$\mathbb R\oplus\bar\rho_r$. Pulling the Thom class back along
the trivial $n$-plane inclusion
\[
 \mathbb R^n\longrightarrow n\rho_r
\]
multiplies it by the Euler class of $n\bar\rho_r$, namely $d_0^n$,
and leaves the trivial Thom shift $\Sigma^n$. Multiplication by
$d_0^n$ is injective in the polynomial algebra. All maps used here
are stable cohomology maps; hence the resulting identification is
the Steenrod module $\Sigma^nd_0^nD_k$ with its inherited action.
In particular this also proves stability of that ideal under $A$.
Weight one is immediate from $D_1(\Sigma^nS)=\Sigma^nS$.

For compatibility with suspension, the Thom map in
Lemma~\ref{lem:weight-suspension} has complement $\bar\rho_r$.
Its pullback sends a Thom coefficient $f$ to $d_0f$. The level-$n$
identification then sends this coefficient to
$d_0^n(d_0f)=d_0^{n+1}f$, which is its image under the level-$(n+1)$
identification. Thus, after the common suspension shift, the map is
\[
 M_k(n+1)=\Sigma^{n+1}d_0^{n+1}D_k
 \longrightarrow
 \Sigma M_k(n)=\Sigma^{n+1}d_0^nD_k,
\]
given by inclusion. This proves \eqref{eq:module-suspension}, with
all degree shifts as asserted.
\end{proof}

\subsection{The weight filtration at level zero}

At positive levels the structured splitting gives a direct sum of
weights. At level zero we instead construct compatible finite weight
projections and show that they determine the first nonzero weight.
This requires a truncation of the free algebra.

For $N\geq0$, write $\mathbb P_{\leq N}(T)$ for the algebra obtained by
discarding weights greater than $N$. Its underlying spectrum is
$\bigvee_{0\leq r\leq N}D_rT$, and an operation whose output weight exceeds
$N$ is zero. More formally, one first truncates in nonnegatively graded
spectra, with the truncated tensor product. Extend the result by zero
to the full category of nonnegatively graded spectra, using the
induced lax monoidal structure, and then take its direct sum.
Nonnegative grading ensures that an operation killed
by the truncation cannot subsequently return to a retained weight. This
gives a commutative ring spectrum and a natural algebra projection
$\mathbb P(T)\to\mathbb P_{\leq N}(T)$.

The particle monoid at level zero is
$C(S^0)=\coprod_{r\geq0}B\Sigma_r$. It is essential to distinguish
its particle generator from the reduced generator of
Proposition~\ref{prop:free-ring}. Denote the former by $t$ and the
latter by $u$, so that $t=1+u$. The symbol $1$ here is the additive
monoid unit, represented in homology by the empty-particle component
point; it is not the point in component $-1$ of $QS^0$.

\begin{lemma}\label{lem:zero-weights}
There is an equivalence of augmented commutative ring spectra
\begin{equation}\label{eq:zero-localization}
 R:=\Sigma^\infty_+QS^0
   \simeq\mathbb P(S\{u\})[(1+u)^{-1}],\qquad t=1+u,
\end{equation}
and compatible augmented commutative ring maps
\begin{equation}\label{eq:weight-projections}
 R\longrightarrow\mathbb P_{\leq N}(S\{u\}),\qquad N\geq0.
\end{equation}
The induced homology maps jointly detect all classes. Consequently a
nonzero positive-dimensional class has a well-defined smallest
nonzero weight, called its \emph{leading weight}. Multiplication by
a power of $t$ does not change its leading component.

Let $z\in H_dQ_0S^0$, $d>0$, be primitive and annihilated by positive
Steenrod operations. Its leading component $z_r$ is primitive and
$A$-annihilated for the linear coproduct calculation of
Lemma~\ref{lem:weight-modules} at $n=0$. In particular
$r=2^k$ for some $k\geq1$. The weight-$r$ component of
$\sigma_*z$ is the image of $z_r$ under
\eqref{eq:weight-suspension}; its functional on $M_k(1)$ is the
restriction, shifted by one, of the functional on $D_k$ defined by
$z_r$. If $d$ is odd, this component is nonzero.
\end{lemma}

\begin{proof}
Under Proposition~\ref{prop:free-ring}, the particle is the ring
element $t=1+u$. Group completion of $C(S^0)$ inverts this particle
component. For a commutative ring spectrum $A'$, a map from its
suspension ring is a map from the additive $E_\infty$ particle monoid
to the multiplicative $E_\infty$ space of $A'$. Such a map extends
over group completion precisely when the particle maps to a unit.
Equivalently, the resulting ring map inverts $1+u$. This is the
universal property of the localization in
\eqref{eq:zero-localization}, including its mapping spaces. The
augmentation sends $u$ to zero and $t$ to one, so the equivalence is
augmented. The underlying module of localization at a degree-zero
element is its multiplication telescope. Therefore ordinary
homology satisfies
\[
 H_*R=H_*\mathbb P(S\{u\})[(1+u)^{-1}].
\]

In $\mathbb P_{\leq N}(S\{u\})$, the element $u$ is nilpotent in
$\pi_0$. Thus $1+u$ is a unit, with inverse
$\sum_{j=0}^N(-u)^j$ in $\pi_0$, and the truncation extends across
localization to give \eqref{eq:weight-projections}. Their
compatibility follows from the localization universal property.

Every class of $H_*R$ can be multiplied by a sufficiently large
power of $t$ to become a class in the free homology algebra
$H_*\mathbb P(S\{u\})$. This algebra is polynomial, as used in
Lemma~\ref{lem:weight-modules}, and is a direct sum of weights.
Hence such a numerator has only finitely many nonzero weights.
If its smallest nonzero component has weight $r$, multiplying by
$(1+u)^a$ leaves that component unchanged: the weight-zero term is
one and every other term has positive weight. Conversely, in a
finite truncation the inverse of $1+u$ has weight-zero term one
and only positive weights otherwise. Thus a nonzero numerator
remains nonzero in every truncation retaining its leading weight.
This proves joint detection, independence of the leading component
from the chosen numerator, and invariance under powers of $t$.
The inclusion $H_*Q_0S^0\to H_*QS^0$ is injective, since it is the
inclusion of one component, so these conclusions apply to the base
component as well.

Before localization, the actual space diagonal on the reduced
generator is
\begin{equation}\label{eq:zero-diagonal}
 \Delta u=u\otimes1+1\otimes u+u\otimes u.
\end{equation}
Consequently, a weight-$r$ class maps only to terms of total weight
at least $r$, and its terms of total weight $r$ are given by the
linear coproduct. This follows by applying extended powers to
\eqref{eq:zero-diagonal} and then multiplying. It remains true
after localization: the diagonal of $t$ is $t\otimes t$, whose
inverse has constant term one and only positive total weights
otherwise. These statements may all be made in finite truncations
in the two target factors, so no convergence of an infinite sum
is required.

The kernels of the projections through weight $r-1$ define a decreasing
weight filtration. The preceding calculation identifies its associated
graded coproduct with the linear coproduct.
Since $z$ is primitive, comparison in its leading total weight gives
\[
 \Delta_{\mathrm{lin}}z_r=z_r\otimes1+1\otimes z_r.
\]
The weight projections are maps of spectra, so they commute with
the stable Steenrod action. Thus $z_r$ is also $A$-annihilated.
Lemma~\ref{lem:weight-modules} makes $r$ a power of two; weight
one has no positive-dimensional homology, so $r=2^k$ with $k\geq1$.

For the suspension assertion, choose $a\geq0$ such that $t^az$
has a finite numerator in the particle monoid homology. The
translation rule in Lemma~\ref{lem:weight-suspension} gives
\[
 \sigma_*(t^az)=\sigma_*z.
\]
The numerator has the same leading component $z_r$ and no smaller
weights. Apply the weight-preserving family of
Lemma~\ref{lem:weight-suspension} before group completion.
It follows that the weight-$r$ component of the suspension is
exactly the image of $z_r$ under
\eqref{eq:weight-suspension}. By
\eqref{eq:module-suspension}, the corresponding functional is the
restriction from $D_k$ to $d_0D_k$, with one suspension shift.

Finally, all the generators of $D_k/(d_0)$ have even degrees:
$|d_j|=2^k-2^j$ is even for $j\geq1$ (and for $k=1$ the quotient
is just $\F_2$). A nonzero odd-degree functional on $D_k$ therefore
cannot vanish on $d_0D_k$. This proves that the leading component
survives suspension when $d$ is odd.
\end{proof}

\subsection{The suspension kernel and square roots of primitives}

Let $H^+$ denote the positive-degree part of a connected homology
algebra, and call $(H^+)^2$ its subspace of decomposables.

\begin{lemma}\label{lem:suspension-kernel}
For $n\geq1$, a positive-dimensional class in the kernel of
\[
 \sigma_*:H_*QS^n\longrightarrow H_{*+1}QS^{n+1}
\]
is decomposable. The same assertion holds for
$\sigma_*:H_*Q_0S^0\to H_{*+1}QS^1$. In each case the kernel in
positive dimensions is precisely the subspace of decomposables.
\end{lemma}

\begin{proof}
The product rule of Lemma~\ref{lem:weight-suspension} kills the
decomposables. For the reverse inclusion, use the polynomial
generators in \eqref{eq:free-homology-basis}. Suspension commutes
with upper-index operations. When the bottom class changes from
$g_n$ to $g_{n+1}$, every intermediate word with those same upper
indices increases in degree by one. Thus each lower index of a
nonempty polynomial generator decreases by one. The original
indices satisfy $0<j_1\leq\cdots\leq j_k$, so the resulting
indices satisfy
\[
 0\leq j_1-1\leq\cdots\leq j_k-1.
\]
These are distinct nonzero primitive words in the target. The
empty-word generator $g_n$ suspends to $g_{n+1}$, also nonzero.
Hence suspension is injective on the span of the polynomial
generators modulo decomposables.

At level zero, localization at the particle gives
\[
 H_*QS^0=\F_2\bigl[t,t^{-1},
 Q_{j_1}\cdots Q_{j_k}t:
 k\geq1,\ 0<j_1\leq\cdots\leq j_k\bigr].
\]
A word of length $k$ has component degree $2^k$. Replacing each
such generator by its product with $t^{-2^k}$ identifies this algebra
with the Laurent polynomial algebra on $t$ over a polynomial algebra
concentrated in component zero. Thus $H_*Q_0S^0$ is polynomial,
with generators the translates to component zero
of the nonempty positive-excess particle generators
$Q_{j_1}\cdots Q_{j_k}t$ with
$0<j_1\leq\cdots\leq j_k$. A word of length $k$ is translated
by $t^{-2^k}$. Suspension ignores that translation and commutes
with its upper-index operations, giving the same distinct
nonzero words on $g_1$ with all lower indices decreased by one.
One may equivalently start with the reduced generator $u=t-1$:
positive operations on the additive monoid unit vanish, so the
positive operations on $u$ and $t$ agree. This proves injectivity
on indecomposables at level zero as well.

The injectivity here is into the full homology of the target.
Some of the resulting words have initial lower index zero and
are themselves squares; they must not be discarded by projecting
the target to indecomposables. Their independence as primitive
words is exactly what gives the asserted kernel description.
\end{proof}

\begin{lemma}\label{lem:primitive-roots}
Let $H$ be a connected graded commutative Hopf algebra over $\F_2$
which is polynomial as an algebra on positive-degree generators.
Every decomposable primitive of $H$ is a square, and its unique
square root is primitive. If $H$ also has a Steenrod action
satisfying the homology Cartan formula, an $A$-annihilated square
has an $A$-annihilated root.

These statements apply to $H_*Q_0S^0$ and to $H_*QS^n$ for
$n\geq1$.
\end{lemma}

\begin{proof}
Put $QH=H^+/(H^+)^2$, let $\pi:H\to QH$ send $1$ to zero,
and define
\[
 \partial=(1\otimes\pi)\Delta:H\longrightarrow H\otimes QH.
\]
The identity
$\pi(ab)=\epsilon(a)\pi(b)+\epsilon(b)\pi(a)$ shows that
$\partial$ is a derivation. If $y_i$ are homogeneous polynomial
generators, connectedness and grading give
\[
 \partial y_i=1\otimes[y_i]
  +\text{terms with smaller positive degree in the second factor}.
\]
For a polynomial $f$, the derivation rule yields
\[
 \partial f=\sum_i\frac{\partial f}{\partial y_i}\,\partial y_i.
\]
Only finitely many variables occur in $f$. If $\partial f=0$,
inspect the coefficient of $[y_i]$ for variables of largest
degree occurring in $f$. There is no contribution from a
higher-degree generator, so their ordinary polynomial partial
derivatives vanish. Descending through generator degrees shows
that every $\partial f/\partial y_i$ vanishes. Over $\F_2$, this
means that every exponent in every monomial of $f$ is even;
equivalently, $f=z^2$ for a unique polynomial $z$.

If $f$ is a decomposable primitive, then
$\partial f=1\otimes\pi(f)=0$, so this argument applies.
Moreover,
\[
 \bigl(\Delta z-z\otimes1-1\otimes z\bigr)^2=0.
\]
The tensor product $H\otimes H$ is again a polynomial algebra,
so its Frobenius map is injective. Hence the displayed reduced
coproduct is zero and $z$ is primitive.

Finally, the homology Cartan formula in characteristic two gives
\[
 \Sq_*^{2a}(z^2)=(\Sq_*^a z)^2,
 \qquad \Sq_*^{2a+1}(z^2)=0.
\]
If $z^2$ is $A$-annihilated, injectivity of Frobenius forces
$\Sq_*^a z=0$ for every $a>0$. Since these squares generate the
positive Steenrod action, the root is $A$-annihilated. The
polynomial homology descriptions above verify the algebra
hypotheses for the infinite loop spaces under consideration.
\end{proof}
